\clearpage\fancyhead[L]{\small COMET.}\hypertarget{native-387-lower}{}\label{native-387-lower}
\begin{center}\textbf{COMET.}\par\vspace{4pt}\hrule\vspace{4pt}\fontsize{10}{12}\selectfont Göttingische gelehrte Anzeigen. No. 55. pp. 537\;..\;538. 8 April 1815.\par\vspace{4pt}\hrule\vspace{4pt}\end{center}

\par\indent From 13 to 31 March there were only four clear evenings here when the new comet could be observed. Here, then, are all the observations obtained so far at the local observatory:

\begin{center}\fontsize{9}{11}\selectfont\setlength{\tabcolsep}{3pt}\renewcommand{\arraystretch}{1.15}\begin{tabular}{l|ccc|ccc|ccc}
1815 & \multicolumn{3}{c|}{Mean time} & \multicolumn{3}{c|}{\shortstack{Apparent right\\ascension}} & \multicolumn{3}{c}{\shortstack{Apparent\\declination}} \\

\hline
March 20 & $10^{\mathrm u}$ & $33^{\mathrm m}$ & $6^{\mathrm s}$ & $54^\circ$ & $7'$ & $1''$ & $39^\circ$ & $7'$ & $47''$ N. \\
21 & 10 & 11 & 37 & 54 & 33 & 21 & 39 & 36 & 57 \\
21 & 10 & 57 & 0 & 54 & 34 & 21 & 39 & 38 & 8 \\
25 & 9 & 46 & 38 & 56 & 28 & 50 & 41 & 38 & 5 \\
30 & 9 & 50 & 57 & 59 & 13 & 3 & 44 & 10 & 27 \\
\end{tabular}\end{center}

\par\indent The second observation on 21 March is by Mr. \textsc{Enke}, whose skill in practical astronomy has often been evident in these pages. No further observations from elsewhere have yet arrived.

\Needspace{19\baselineskip}\par\indent Professor \textsc{Gauss} made the following preliminary, approximate determination of the comet’s orbit from these admittedly scant observations:

\begin{center}\fontsize{10}{12}\selectfont\setlength{\tabcolsep}{3pt}\renewcommand{\arraystretch}{1.15}\begin{tabular}{l}
Passage through perihelion\\
\quad Mean time in Göttingen, 24 April 1815. $16^{\mathrm u}\;37^{\mathrm m}\;34^{\mathrm s}$\\
\end{tabular}\end{center}
\begin{center}\fontsize{10}{12}\selectfont\setlength{\tabcolsep}{3pt}\renewcommand{\arraystretch}{1.15}\begin{tabular}{lcr}
Longitude of perihelion & $.\;.\;.\;.\;.$ & $246^\circ\;7'\;2''$\\
Ascending node & $.\;.\;.\;.\;.\;.$ & $82\;43\;4$\\
Inclination of orbit & $.\;.\;.\;.\;.\;.\;.$ & $45^\circ\;8'\;55''$\\
Least distance from the Sun & $.\;.$ & $1.24738$\\
\end{tabular}\end{center}
\noindent Prograde motion.

\SourceLeaf{388}\par\indent From this determination it can already be predicted that the comet will grow somewhat brighter until near the end of April. The increase is very slight, however. The comet will remain invisible to the naked eye and will therefore be of interest \emph{only to astronomers}. It will pass through Auriga and Lynx into the Great Bear and will disappear from view there in June as its light continues to diminish.
